\section{Introduction}
\label{sec:intro}

Deductive qualitative coding applies a fixed codebook to text: each segment of an interview, dialogue or survey answer is
read, compared with the code definitions and labelled, often twice to establish reliability \citep{krippendorff2019content}.
Large language models (LLMs) can apply codebooks with useful accuracy
\citep{xiao2023supporting,chew2023llm,dunivin2024scalable,ziems2024can,gilardi2023chatgpt,matveyenko2025blended}, but
accuracy alone does not tell a researcher which model labels to trust. The practical promise is \emph{triage}: accept the
labels a model is confident about and send the rest to a person.

Consider a researcher with 1,500 counselling utterances to code and time to check 400 by hand. A model labels every
utterance and attaches a confidence; a threshold decides which labels are accepted. Three properties of the confidence then
matter, and they are easy to conflate. \emph{Error ranking}: do the model's mistakes get lower confidence than its correct
labels? \emph{Calibration}: does a stated 90\% mean 90\% correct? \emph{Selective risk and coverage}: at the chosen threshold,
what share of labels is accepted, and what share of those is wrong? A well-calibrated model can rank its errors poorly, and a
threshold chosen on 400 hand-coded utterances to give 5\% error can give more than 5\% on the remaining 1,100. Our question is
whether model confidence can select low-error decisions reliably enough for this workflow.

A new class of products is marketed for exactly this. We call them \emph{decision models}: an interface that takes typed
questions (yes/no, multiple choice) with criteria and returns a probability for each option, without generating text.
TypeSafe's Jev \citep{typesafe2026jev} was followed within weeks by OpenAI's Decisions API \citep{openai2026decisions},
Cloudflare's Clef \citep{cloudflare2026clef}, Perplexity's pplx-decider \citep{perplexity2026decider} and the open Kev
models \citep{palmer2026kev}. The interface is new; the underlying idea, probabilistic classification with abstention, is
not \citep{geifman2017selective}. Whether these probabilities support triage on real codebooks, with human disagreement,
dialogue context and rare codes, is an open question. We study it for the classification step of deductive coding on short
segments (comments, utterances, sentences), not for inductive coding or long-form interviews.

We ask three research questions, each with pre-specified contrasts:
\begin{itemize}
  \item \textbf{RQ1 (triage).} Do decision models' probabilities rank their own coding errors better than LLM confidence
  signals, and what share of decisions can be accepted at $\le$5\% error?
  \item \textbf{RQ2 (disagreement).} Is model uncertainty high where human coders disagree, and how does model agreement with
  annotators compare with annotators' agreement with each other?
  \item \textbf{RQ3 (robustness).} How much do decisions change under surface changes to the codebook and request, and under
  removal of context?
\end{itemize}
Six experimental questions explain the answers: \textbf{EQ1}, does calibration error predict triage? \textbf{EQ2}, do
threshold rules keep their error target, and what does a finite-sample guarantee cost? \textbf{EQ3}, what limits triage,
ranking or ties at the top? \textbf{EQ4}, do related 9B decision models outperform a 9B instruct model's own confidence?
\textbf{EQ5}, which configurations give the most guaranteed automation per dollar? \textbf{EQ6}, are the results artefacts
of nondeterminism, prompt choice, memorization or label noise?

The short answer to the title is: on Reddit comments, yes, but much of what is accepted is the easy call that a code does
not apply; on dialogue, with whole transcripts hand-coded, our guaranteed rule certified nothing. Two facts drive this. Coding
decisions come in clusters (eight codes of one comment, many utterances of one transcript), so an error guarantee must respect
how the hand-coded sample was drawn: a rule that treats the decisions of hand-coded transcripts as independent fails, and with
whole transcripts of very unequal length as sampling units, our valid procedure cannot certify even a perfect coder at our
calibration sizes; for three of the four dialogue tasks its necessary bound exceeds the available corpus. Sampling single
utterances instead removes that obstacle, but at our budget the coders' confident dialogue decisions still give too little
evidence of $\le$5\% error. And most decisions in a multi-label codebook are negatives, which a baseline that never reads the text accepts
at low error.

\paragraph{Relation to concurrent work.} Two studies released in the weeks before ours also evaluate decision models:
\citet{ibrahim2026decision} compare Jev and small open decision models with 19 LLMs on 15 computational-social-science labelling
tasks (accuracy, calibration, coverage at fixed thresholds, cost and cascades), and TypedBench \citep{sharma2026typedbench} uses
synthetic, rule-labelled tasks. Where the designs overlap, our findings agree: Jev trails the strongest LLMs in accuracy, its raw
calibration beats most LLM signals but not the frontier models, and cascades save cost. Risk-controlled triage is not new
either: HyPAC \citep{zeng2026hypac} routes items between models and humans with a PAC guarantee on annotation error, and
\citet{zhao2025confidence} triage qualitative codes by confidence and inter-model diversity. Our contribution is an evaluation,
not a new method.

\paragraph{Contributions.}
\begin{enumerate}
  \item \textbf{A multi-vendor evaluation.} Seven decision models and twelve LLM confidence signals on five human-coded
  tasks, under a prospectively specified plan with a deviations log, with every model output released (\S\ref{sec:setup}).
  \item \textbf{Ranking, not calibration error, separates coders.} Raw error ranking varies widely across decision models,
  low calibration error does not predict it, and per-code recalibration narrows most differences (\S\ref{sec:rq1},
  \S\ref{sec:eq1}--\S\ref{sec:eq3}).
  \item \textbf{Guarantees that respect the sampling design.} The usual rule overshoots its target, and a guarantee that treats
  decisions as independent fails. A corpus-level rule adapted from betting tests for sampling without replacement certified
  most comment decisions and, with whole transcripts hand-coded, no dialogue decisions; we separate what the sampling design,
  the procedure and the coders each contribute to that cost (\S\ref{sec:rules}, \S\ref{sec:eq2}).
  \item \textbf{What automation means for coding.} A prevalence baseline, recall of applying codes, the share of segments still
  read by a person and policy-specific cost give a more honest picture than accepted share alone (\S\ref{sec:eq5},
  \S\ref{sec:discussion}).
\end{enumerate}
